\documentclass[10pt,a4paper]{amsart}
\usepackage[margin=19mm]{geometry}
\usepackage{amsmath,amssymb,mathtools}
\usepackage[scr=rsfs,cal=euler]{mathalfa}
\usepackage{xfrac}
\usepackage[unicode,hidelinks,bookmarks=false]{hyperref}
\newcommand{\ie}{i.e.\ }
\newcommand{\cf}{cf.\ }
\newcommand{\resp}{respectively\ }
\newcommand{\Subsection}{\subsection}
\providecommand{\nobreakdash}{\nobreak\hbox{-}}
\setlength{\emergencystretch}{2em}
\multlinegap0pt
\usepackage{amssymb}
\usepackage{tcolorbox}
\usepackage{enumerate}
\newtheorem{prop}{Proposition}
\DeclareMathOperator{\bbz}{\mathbb{Z}}
\title{Revisiting fully residually free Demushkin groups}
\author{Henrique Souza, Pavel Zalesskii}
\date{Source: arXiv:2603.16814v1 (CC BY 4.0)}
\begin{document}
\maketitle

\noindent\emph{Source and licence.} Revisiting fully residually free Demushkin groups. Version arXiv:2603.16814v1 (CC BY 4.0), distributed by arXiv under the Creative Commons Attribution 4.0 International licence, http://creativecommons.org/licenses/by/4.0/, as recorded in the arXiv licence metadata for this exact version and shown on the abstract page https://arxiv.org/abs/2603.16814v1. Full licence notice: \texttt{licenses/arxiv-2603.16814-CC-BY-4.0.txt}.\par\smallskip

\noindent\textit{Editorial scope.} Counted unit: the article's prop prop-lyndon (source0.tex lines 263--264). The article's complete original proof of it is delivered with it (source0.tex lines 265--275, one \texttt{proof} environment of 11 lines). No mathematics is written, repaired or re-derived: the statement and the proof are the article's own lines, in the article's own notation, and no retained line is substituted or re-flowed.  No further source-local context is needed: every cross-reference of every retained line resolves inside the two delivered spans.  Nothing was omitted from any retained span, so the package carries no omission record.  No retained line contains a citation, so the wrapper carries no bibliography and no bibliography entry.  No retained line contains a figure environment, an image-inclusion command or drawing code, so no image is delivered and the package declares no asset dependency.  Editorial formatting only, in addition: an AMS \texttt{amsart} preamble that restates the article's own declarations and macros used by the retained lines, namely the environments \texttt{prop} on separate counters, together with the standard packages those lines need.  The retained environments print in this document as Proposition 1 in order of appearance, whereas the article prints the same items with the numbering of its own full document.  The complete native source of the article as distributed in the arXiv e-print for this exact version is bundled unchanged as \texttt{sources/196580/source0.tex}.
\begin{prop}\label{prop-lyndon} Let \(F\) be a free pro-\(2\) group and \(a,b,c \in F\) be three elements satisfying \(a^2b^2c^2 = 1\). Then, \(\langle a,b,c\rangle\) is abelian.
\end{prop}

\begin{proof} Let \(H\) be the closed subgroup of \(F\) generated by \(a\), \(b\) and \(c\). It is free pro-\(p\), and any set whose image in \(H/H^2\) is a basis is a free generating set for \(H\). If \(\{a,b,c\}\) were a free generating set for \(H\), we could obtain a homomorphism \(H \to \mathrm{GL}_3(\bbz/4\bbz)\) given by \[a \mapsto A = \begin{pmatrix}
    1 & 1 & 0\\
    0 & 1 & 0\\
    0 & 0 & 1\\
\end{pmatrix}\,,\quad b \mapsto B = \begin{pmatrix}
    1 & 0 & 0\\0 & 1 & 1\\
    0 & 0 & 1
\end{pmatrix}\,,\quad c \mapsto C = \begin{pmatrix}
    1 & 0 & 1\\0 & 1  & 0\\0 & 0 & 1
\end{pmatrix}\,.\] However, one readily checks that the matrix \(A^2B^2C^2\) is not the identity matrix, yielding a contradiction. Hence, there is a linear dependence between the \(a,b,c\) modulo \(H^2\). Without loss of generality, we can then assume that \(H\) is generated by \(a\) and \(b\). If \(\{a,b\}\) were a free basis of \(H\), we could also obtain a homomorphism \(f\colon H \to \mathrm{GL}_3(\bbz/4\bbz)\) by mapping \(a\) to \(A\) and \(b\) to \(B\). Since the matrix \(A^2B^2\) has no square root in \(f(H)\), this also gives a contradiction. Therefore, \(a\) and \(b\) are linearly dependent modulo \(H^2\) and \(H\) must be procyclic.
\end{proof}

\end{document}
